% ============================================================
\chapter{Lebesgue Measure on $\R^n$}
\label{ch:lebesgue}
% ============================================================

We finally arrive at the object that all the preceding machinery has been preparing: the \emph{Lebesgue measure}. It is the measure that gives rigorous meaning to the notions of length, area, and volume, where Riemann's construction fell short. Henri Lebesgue, in his 1902 thesis, showed that one can measure far more pathological sets than simple unions of intervals---and that the resulting integral possesses limit-passage properties incomparably superior to Riemann's. The construction we present here applies Carath\'eodory's theorem to the volume content of rectangles, and produces the most fundamental measure in all of analysis.

\section{Construction}

Lebesgue measure is the unique measure on $\R^n$ that extends the notion of volume
of rectangles. We construct it via Carath\'eodory's theorem.

\begin{definition}[Rectangle and volume]
A \emph{rectangle} (or \emph{box}) in $\R^n$ is a set of the form
$I = \prod_{k=1}^n [a_k, b_k)$ with $a_k \leq b_k$. Its \emph{volume} is:
\[
\operatorname{vol}(I) = \prod_{k=1}^n (b_k - a_k).
\]
\end{definition}

\begin{definition}[Algebra of disjoint rectangles]
Let $\mathcal{A}_0$ denote the collection of finite disjoint unions of half-open
rectangles in $\R^n$. This is an algebra.
\end{definition}

\begin{proposition}
The function $\mu_0 : \mathcal{A}_0 \to [0, +\infty]$ defined by
$\mu_0\!\left(\bigsqcup_{k=1}^m I_k\right) = \sum_{k=1}^m \operatorname{vol}(I_k)$
is a $\sigma$-finite pre-measure on $\mathcal{A}_0$.
\end{proposition}

\begin{proof}
Finite additivity is immediate by definition. The $\sigma$-additivity follows from
continuity at $\varnothing$ (Proposition~\ref{prop:sigma_add_criterion}). Let
$(A_m)$ be a decreasing sequence in $\mathcal{A}_0$ with $A_m \downarrow
\varnothing$. We show $\mu_0(A_m) \to 0$ by a compactness argument.

Suppose for contradiction that $\mu_0(A_m) \geq \alpha > 0$ for all $m$. Each
$A_m$ is a finite union of rectangles. By slightly shrinking each rectangle
$[a_k, b_k)$ to $[a_k + \varepsilon, b_k - \varepsilon]$, we obtain a compact set
$K_m \subset A_m$ with $\mu_0(A_m) - \mu_0(K_m)$ small. We can ensure
$\mu_0(K_m) \geq \alpha/2 > 0$, so $K_m \neq \varnothing$. Since
$K_1 \supset K_2 \supset \ldots$ are nested non-empty compact sets,
$\bigcap_m K_m \neq \varnothing$, contradicting $\bigcap_m A_m = \varnothing$.

$\sigma$-finiteness follows from $\R^n = \bigcup_{k=1}^{\infty} [-k, k)^n$.
\end{proof}

\begin{theorem}[Lebesgue measure]
\label{thm:lebesgue_measure}
There exists a unique measure $\lambda^n$ on $(\R^n, \mathcal{B}(\R^n))$ such that:
\[
\lambda^n\!\left(\prod_{k=1}^n [a_k, b_k)\right) = \prod_{k=1}^n (b_k - a_k).
\]
This measure is called \emph{Lebesgue measure} on $\R^n$.
\end{theorem}

\begin{proof}
Direct application of the Carath\'eodory extension theorem
(Theorem~\ref{thm:caratheodory_extension}) to the pre-measure $\mu_0$. Uniqueness
follows from $\sigma$-finiteness.
\end{proof}

\begin{remark}
We often write $\lambda = \lambda^1$ for Lebesgue measure on $\R$. The
\emph{complete Lebesgue measure} is the completion of $\lambda^n$ with respect to
the Lebesgue $\sigma$-algebra $\mathcal{L}^n \supset \mathcal{B}(\R^n)$.
\end{remark}

\section{Fundamental properties}

\begin{theorem}[Regularity]
\label{thm:lebesgue_regularity}
For every $A \in \mathcal{B}(\R^n)$:
\begin{enumerate}
    \item (Outer regularity) $\lambda^n(A) = \inf\{\lambda^n(U) :
    U \supset A,\, U \text{ open}\}$.
    \item (Inner regularity) $\lambda^n(A) = \sup\{\lambda^n(K) :
    K \subset A,\, K \text{ compact}\}$.
\end{enumerate}
\end{theorem}

\begin{proof}
(1) Let $\varepsilon > 0$. By definition of the outer measure, there exists a
sequence of rectangles $(I_k)$ with $A \subset \bigcup_k I_k$ and
$\sum_k \operatorname{vol}(I_k) \leq \lambda^n(A) + \varepsilon$. Replacing each
$[a_j, b_j)$ by $(a_j - \delta, b_j + \delta)$ with $\delta$ small enough, we
obtain an open set $U \supset A$ with $\lambda^n(U) \leq \lambda^n(A) +
2\varepsilon$.

(2) For $\lambda^n(A) < \infty$, use the approximation
$A = \bigcup_m (A \cap [-m, m]^n)$ and outer regularity applied to the complement.
\end{proof}

\begin{theorem}[Translation invariance]
\label{thm:translation_invariance}
For every $A \in \mathcal{B}(\R^n)$ and every $x \in \R^n$:
\[
\lambda^n(A + x) = \lambda^n(A), \quad \text{where } A + x = \{a + x : a \in A\}.
\]
\end{theorem}

\begin{proof}
The measure $\nu(A) = \lambda^n(A + x)$ is a measure on $\mathcal{B}(\R^n)$ that
agrees with $\lambda^n$ on rectangles (since volume is translation-invariant). By
uniqueness, $\nu = \lambda^n$.
\end{proof}

\begin{theorem}[Scaling behavior]
For every $A \in \mathcal{B}(\R^n)$ and every $c \in \R \setminus \{0\}$:
\[
\lambda^n(cA) = \abs{c}^n \lambda^n(A).
\]
More generally, if $T : \R^n \to \R^n$ is an invertible linear map:
\[
\lambda^n(T(A)) = \abs{\det T} \cdot \lambda^n(A).
\]
\end{theorem}

\begin{proof}
For scaling, the measure $\nu(A) = \lambda^n(cA)/\abs{c}^n$ agrees with $\lambda^n$
on rectangles and is $\sigma$-finite, so $\nu = \lambda^n$ by uniqueness.

For the general case, reduce to shears and scalings (Gaussian elimination).
\end{proof}

\section{Null sets}

\begin{proposition}
The following sets have Lebesgue measure zero:
\begin{enumerate}
    \item every countable set;
    \item every affine subspace of dimension $< n$ in $\R^n$;
    \item the Cantor set $\mathcal{C} \subset [0,1]$.
\end{enumerate}
\end{proposition}

\begin{proof}
(1) A singleton $\{x\}$ is contained in $[x_1, x_1 + \varepsilon) \times \ldots
\times [x_n, x_n + \varepsilon)$ of volume $\varepsilon^n$, for every
$\varepsilon > 0$. So $\lambda^n(\{x\}) = 0$. By countable subadditivity, every
countable set has measure zero.

(2) By isometry invariance, it suffices to treat the hyperplane
$\R^{n-1} \times \{0\}$. We have $\R^{n-1} \times \{0\} \subset \R^{n-1} \times
[-\varepsilon, \varepsilon)$ for all $\varepsilon > 0$, and the latter has finite
measure tending to zero.

(3) The Cantor set is the intersection $\mathcal{C} = \bigcap_n C_n$ where $C_n$ is
a union of $2^n$ intervals of length $3^{-n}$. So
$\lambda(\mathcal{C}) \leq \lambda(C_n) = (2/3)^n \to 0$.
\end{proof}

\section{Non-measurable sets}

\begin{theorem}[Vitali, 1905]
\label{thm:vitali}
There exists a subset of $[0,1]$ that is not Lebesgue-measurable.
\end{theorem}

\begin{proof}
Define the equivalence relation $x \sim y \iff x - y \in \Q$ on $[0,1)$. By the
axiom of choice, let $V \subset [0,1)$ contain exactly one representative from each
equivalence class.

For $r \in \Q \cap [0,1)$, define $V_r = \{x + r \pmod{1} : x \in V\}$ (translation
modulo 1). The sets $V_r$ are pairwise disjoint and
$[0,1) = \bigsqcup_{r \in \Q \cap [0,1)} V_r$.

If $V$ were measurable, by translation invariance (modulo 1), all $V_r$ would have
the same measure $\alpha = \lambda(V)$. By $\sigma$-additivity:
\[
1 = \lambda([0,1)) = \sum_{r \in \Q \cap [0,1)} \lambda(V_r) =
\sum_{r \in \Q \cap [0,1)} \alpha.
\]
If $\alpha = 0$, the sum equals $0 \neq 1$. If $\alpha > 0$, the sum equals
$+\infty \neq 1$. Contradiction.
\end{proof}

\begin{warning}[title=\textbf{Role of the axiom of choice}]
Vitali's construction uses the axiom of choice. Indeed, Solovay (1970) showed that
it is consistent with ZF (without the axiom of choice) that every subset of $\R$
is Lebesgue-measurable.
\end{warning}

\section{The Cantor set}

\begin{definition}[Cantor set]
The Cantor set $\mathcal{C}$ is defined by:
\[
\mathcal{C} = \bigcap_{n=0}^{\infty} C_n, \quad \text{where }
C_0 = [0,1], \quad C_{n+1} = \frac{C_n}{3} \cup \left(\frac{2}{3} +
\frac{C_n}{3}\right).
\]
Equivalently, $\mathcal{C} = \{x \in [0,1] : x = \sum_{k=1}^{\infty}
a_k 3^{-k},\, a_k \in \{0, 2\}\}$.
\end{definition}

\begin{proposition}[Properties of the Cantor set]
\begin{enumerate}
    \item $\mathcal{C}$ is compact, uncountable, perfect (closed with no isolated
    points), and totally disconnected.
    \item $\lambda(\mathcal{C}) = 0$.
    \item $\mathcal{C}$ is in bijection with $[0,1]$ (via the Cantor function).
    \item $\mathcal{C}$ contains non-Borel subsets (since
    $\abs{\mathcal{P}(\mathcal{C})} = 2^{\mathfrak{c}} > \mathfrak{c} =
    \abs{\mathcal{B}(\R)}$).
\end{enumerate}
\end{proposition}

\section{Characterization of Lebesgue-measurable sets}

\begin{theorem}
A set $A \subset \R^n$ is Lebesgue-measurable if and only if there exist a Borel set
$B$ and a null set $N$ such that $A = B \cup N$. In other words,
$\mathcal{L}^n = \overline{\mathcal{B}(\R^n)}$ (the completion of the Borel
$\sigma$-algebra).
\end{theorem}

\begin{theorem}[Characterization via $G_\delta$ and $F_\sigma$]
Let $A \subset \R^n$ be Lebesgue-measurable with $\lambda^n(A) < \infty$. Then:
\begin{enumerate}
    \item there exists a $G_\delta$ set (countable intersection of open sets)
    $G \supset A$ with $\lambda^n(G \setminus A) = 0$;
    \item there exists an $F_\sigma$ set (countable union of closed sets)
    $F \subset A$ with $\lambda^n(A \setminus F) = 0$.
\end{enumerate}
\end{theorem}

\section{Exercises}

\begin{exercise}
Show that $\lambda(\Q) = 0$ directly from the definition of Lebesgue measure (by
constructing an explicit covering).
\end{exercise}

\begin{exercise}
Let $A \subset \R$ be measurable with $\lambda(A) > 0$. Show that the difference
set $A - A = \{a - b : a, b \in A\}$ contains an interval centered at zero.
(\emph{Steinhaus' theorem.})
\end{exercise}

\begin{exercise}[Generalized Cantor set]
For $\alpha \in (0, 1)$, construct a closed set $\mathcal{C}_\alpha \subset [0, 1]$,
homeomorphic to the Cantor set, with $\lambda(\mathcal{C}_\alpha) = 1 - \alpha$.
\end{exercise}

\begin{exercise}
Let $f : \R \to \R$ be Lipschitz with constant $L$. Show that for every
$A \in \mathcal{B}(\R)$, $\lambda(f(A)) \leq L \cdot \lambda(A)$.
\end{exercise}

\begin{exercise}
Show that Lebesgue measure $\lambda$ is the only (complete) measure on $\R$ that is
translation-invariant and satisfies $\lambda([0,1]) = 1$.
\end{exercise}

\begin{exercise}
Show that there exists a Lebesgue-measurable set that is not Borel.
\emph{Hint}: use the surjection $\mathcal{C} \to [0,1]$ and the fact that
$\mathcal{C}$ has measure zero.
\end{exercise}
